% Page budget: 1.55 pages | Owns: augmentation topology, additional states, encoder, and closed-loop objective.
% WRITING GUIDE
% Purpose: Define the final augmentation architecture and explain how it is trained inside the known feedback loop.
% Start here: Current implementation decisions D-129, D-140, D-141, D-180, D-220, D-233, and D-237 in docs/decisions.md.
% Supporting material: Quinten's 18 September outline feedback, the Meeting-audit synthesis, Thesis-writeup/Documentation/TRAINING-DESIGN.md, and docs/writeup figures.
% Mandatory papers: Read Hoekstra's 2025 LFR augmentation paper, the 2026 first-principles augmentation paper, the encoder-initialization paper, and Kessels' Chapter 5 before writing this section.
% Research-plan anchor: Preserve Aspect 2's dynamic parallel augmentation objective; document the departure from open-loop training and earlier state routing.
% Current decisions: Separate physical and additional states, keep the controller outside the model, score the complete training window without burn-in, and use one network for the physical correction and added-state update.
% Choices to justify: Final reported routing, added-state dimension, ANN width and depth, zero initialization, encoder history, closed-loop objective, rollout length, full-window scoring, and validation horizon. Use the configuration of the reported thesis run, not a superseded routing experiment.
% Horizon check: Treat encoder history, scored training window, joint-estimation horizon, and free-run validation horizon separately; relate candidates to relevant stable modes and test sensitivity to slower dynamics and free-integrator accumulation.
% Implementation audit: Read scripts/gantry/gantry_interconnect_dynamic.py, gantry_dynamic/{model,config,controller,data,training}.py, and model_augmentation/fit_systems/{interconnect,closed_loop,pre_encoder}.py.
% Reference check: Verify sources for parallel LFR augmentation, encoder initialization, closed-loop state extension, and any claimed architectural precedent against the original paper or thesis.
% Cautions: Earlier routing plans are superseded; distinguish the truth-only hidden absorber from learned states; closed-loop fit alone does not prove autonomous plant recovery.
% Figures: Absorber schematic, author's updated architecture, and thesis-owned common-reference loops. No residual panel or horizon figure.
% Section is finished when: Every signal has a defined frame and timing, the RK4 boundary is explicit, and the described architecture matches the final code.
%
% SOURCE MAP
% Hoekstra et al. 2025 EJC: dynamic-parallel topology, additional states, encoder-based truncated simulation.
% Hoekstra et al. 2026 Automatica submission: extended LFR formulation, baseline-preserving model/encoder initialization, joint identification context.
% Hoekstra et al. 2026 encoder paper: reconstructability-based W^b initialization; it does not identify W^a or test dynamic augmentation.
% Kessels 2025 thesis: direct closed-loop rollout, controller equations, and reconstruction of the machine controller state for each window.
% Drenth 2025 thesis: LPV-specific augmentation precedent; not the unrestricted MLP structure used here.
% This project: gantry specialization, residual-loop implementation, exact routing, RK4 boundary, and verification.
\section{Dynamic LPV-LFR Augmentation}\label{sec:augmentation}

The baseline of Section~\ref{sec:lfr} describes the rigid-body dynamics of
the gantry, but not every effect of the machine.
Garc\'ia-Herreros et al.\ neglect the resonance of the supporting base, the
vibration of the cross-arm on its flexible plates, the detent forces of the
linear motors, and the variation of friction along the
stroke~\cite[Sec.~2.4]{garcia2013model}, and our baseline \eqref{eq:eom}
also leaves out the Coulomb friction that their model contains.
This section augments the baseline with a learned component that has states
of its own, and estimates the ten physical combinations
$\theta_{\mathrm{base}}$ of Section~\ref{sec:params}, the network parameters
$\theta_{\mathrm{aug}}$ and the encoder parameters $\eta$ jointly from records
measured under feedback.
Section~\ref{sec:aug_structure} defines the augmented model and its behaviour
at initialisation, Section~\ref{sec:aug_encoder} the encoder that supplies the
initial state of each training window, and Section~\ref{sec:aug_closed_loop}
the closed-loop rollout of each window and the objective minimised over the
windows.

\subsection{Augmented model}\label{sec:aug_structure}
An omitted vibration mode adds model order: the machine response contains
poles that the six states of the baseline cannot produce.
Refitting $\theta_{\mathrm{base}}$ cannot add these poles, and neither can a
static correction of the baseline's state update, because both only change
the dynamics of the existing states; additional states
can~\cite[Sec.~2]{hoekstra2025lfr}.
We therefore adopt the dynamic-parallel structure
of~\cite[Table~1]{hoekstra2026lfr}, in which a learned component with states
of its own acts in parallel to the baseline.
Drenth formulates this kind of augmentation, additional states included, for
a self-scheduled LPV-LFR baseline~\cite[Sec.~5.2]{drenth2025thesis}; here the
baseline is the gantry realisation of Section~\ref{sec:lfr}.
With sample index $k$ and sample period $T_s$, the augmented model is
\begin{equation}
\begin{aligned}
 \tilde x_{k+1}
 &=f_{\mathrm{base}}(\theta_{\mathrm{base}},\tilde x_k,u_k)
   +f_{\mathrm{aug}}(\theta_{\mathrm{aug}},\hat x_k,u_k),\\
 \bar x_{k+1}
 &=g_{\mathrm{aug}}(\theta_{\mathrm{aug}},\hat x_k,u_k),\\
 \hat y_k&=C_y\tilde x_k ,
\end{aligned}
\label{eq:aug_transition}
\end{equation}
where $\tilde x=\col(q,\dot q)\in\R^6$ is the physical state of
Section~\ref{sec:lfr}, $\bar x\in\R^{n_{x_a}}$ collects the additional
states, $\hat x=\col(\tilde x,\bar x)$, $u_k=P\,u_{\mathrm{act},k}$ is the
generalised force of \eqref{eq:Ptransform}, and $\hat y_k$ is the prediction
of the measured positions $y_k=q_{\mathrm{act},k}$.
The baseline transition $f_{\mathrm{base}}$ keeps its physical parameters,
the learned function $f_{\mathrm{aug}}$ corrects the physical states in
parallel to it, and $g_{\mathrm{aug}}$ gives the next additional state.

The baseline transition $f_{\mathrm{base}}$ is one step of the classical
fourth-order Runge-Kutta (RK4) method applied to \eqref{eq:eom} over $T_s$,
with $u_k$ held constant over the step, as the actuator force is held between
samples.
Each of the four stages evaluates the LFR of Section~\ref{sec:lfr} at its own
intermediate state, and therefore at its own payload position $Y$, so the
baseline remains self-scheduled within the step instead of being frozen at
$Y_k$.
The model works in normalised coordinates.
The actuator forces, the measured positions and the physical state are
scaled to zero mean and unit standard deviation over the training records,
and $f_{\mathrm{base}}$ is wrapped by these scalings
as in~\cite[Sec.~5.3, Eq.~(28)]{hoekstra2026lfr}, which changes the
coordinates of the baseline but not its dynamics.
Where~\cite{hoekstra2026lfr} takes the state scale from a simulation of the
baseline, we take it from the training records: the measured positions mapped to $q$
by $P^{-\top}$, and their first differences.
This scaling also defines the inner product of Section~\ref{sec:obc}.
The implementation stores $u_{\mathrm{act}}$ and applies $P$ inside
$f_{\mathrm{base}}$; because $P$ is constant and invertible, this is the same
model.

The two learned functions are the outputs of one multilayer perceptron (MLP)
$\phi_{\mathrm{aug}}$ with tanh hidden layers,
\begin{equation}
 \begin{bmatrix}
  f_{\mathrm{aug}}(\theta_{\mathrm{aug}},\hat x_k,u_k)\\
  g_{\mathrm{aug}}(\theta_{\mathrm{aug}},\hat x_k,u_k)
 \end{bmatrix}
 =\phi_{\mathrm{aug}}\bigl(\theta_{\mathrm{aug}},\col(\hat x_k,u_k)\bigr),
 \label{eq:aug_mlp}
\end{equation}
where $\theta_{\mathrm{aug}}$ collects its weights and biases.
The first six outputs form $f_{\mathrm{aug}}$ and the last $n_{x_a}$ outputs
form $g_{\mathrm{aug}}$.
The additional states therefore have no dynamics outside the network: their
transition and their effect on the physical states are both blocks of the
Jacobian of $\phi_{\mathrm{aug}}$.
The correction $f_{\mathrm{aug}}$ writes every physical row, the $X$ and $Y$
rows included.
The effects listed above act on the $X$ and $Y$ motion, friction at the
actuators being one example, so a correction confined to the stiff $\Theta$
rows could not represent them.
Because the network adds its output to the result of the RK4 step,
$f_{\mathrm{aug}}$ is a discrete-time correction of the sampled transition,
whereas $f_{\mathrm{base}}$ integrates the continuous-time physics.

\begin{figure}[t]
 \centering
 \resizebox{\columnwidth}{!}{\input{figures/tex/augmentation_structure}}
 \caption{Augmented model \eqref{eq:aug_transition}. The physics step
 $f_{\mathrm{base}}$ and the network $\phi_{\mathrm{aug}}$ act in parallel on
 $\hat x_k$; $\phi_{\mathrm{aug}}$ supplies $f_{\mathrm{aug}}$ and
 $g_{\mathrm{aug}}$, and the output map is not learned
 ($h_{\mathrm{aug}}=0$). The encoder $\psi$ sets $\hat x$ at each window
 start. During training, $u_k$ is the closed-loop input $\hat u_k$ of
 Section~\ref{sec:aug_closed_loop}.}
 \label{fig:aug_structure}
\end{figure}

The output map $C_y$ is the kinematic map
$\begin{bmatrix}P^\top & 0\end{bmatrix}$ of \eqref{eq:Ptransform} in
normalised coordinates, so the measured positions are read from the physical
state alone.
We do not augment the output equation, because the measured positions are a
known linear function of the physical state, the case in which Kessels also
leaves the output equation unaugmented~\cite[Remark~5.3]{kessels2025ai}.
Within one step, no learned function feeds the input of another, and
$f_{\mathrm{base}}$ feeds none of them, so the interconnection graph is
acyclic.
The RK4 step of \eqref{eq:eom} is smooth wherever $M(Y)$ is invertible, which
the parameterisation of Section~\ref{sec:params} guarantees, and the tanh
network is smooth everywhere.
Both components are therefore twice differentiable, and the model is well
posed by~\cite[Thm.~9]{hoekstra2026lfr}.
The additional states have no assigned physical meaning; they reach the
output only through $f_{\mathrm{aug}}$, one step after they are formed.

Training starts from the baseline.
The output layer of $\phi_{\mathrm{aug}}$ is initialised at zero,
\begin{equation}
 W_L=0,\quad b_L=0
 \quad\Longrightarrow\quad
 f_{\mathrm{aug}}=0,\quad g_{\mathrm{aug}}=0 ,
 \label{eq:aug_zero_init}
\end{equation}
where $W_L$ and $b_L$ are the weight and bias of the output layer, while the
hidden layers are initialised randomly.
The augmented model then reproduces the baseline from the same physical
initial state, which is the behaviour at initialisation required
in~\cite[Eq.~(29)]{hoekstra2026lfr}, and the additional states are zero
after the first step.
The rows of $W_L$ that produce $f_{\mathrm{aug}}$ receive a gradient at once,
because the random hidden layers give them nonzero inputs.
The rows that produce $g_{\mathrm{aug}}$ do not: $\bar x$ reaches the output
only through $f_{\mathrm{aug}}$, whose dependence on $\bar x$ is zero at
initialisation, so these rows receive a gradient only after training has
given $f_{\mathrm{aug}}$ such a dependence.
From this start, $f_{\mathrm{aug}}$ can also learn changes that a different
$\theta_{\mathrm{base}}$ would produce, so the data do not fix the split
between the two components~\cite[Sec.~5.2]{hoekstra2026lfr}.
Section~\ref{sec:obc} constrains this split.
The number of additional states and the network size are given in
Section~\ref{sec:setup}.

\subsection{Encoder and initial state}\label{sec:aug_encoder}
Each training window simulates \eqref{eq:aug_transition} from an initial
state, but only the positions are measured, so the velocities and the
additional states at the window start are unknown.
Treating them as free variables would add parameters with every
window~\cite[Sec.~1]{hoekstra2026encoder}.
Following Beintema et al.~\cite{beintema2023subnet}, an encoder instead
estimates the initial state from the input and output history,
\begin{equation}
 \hat x_{\tau|\tau}
 =\psi_\eta(\xi_\tau),\qquad
 \xi_\tau=\col\bigl(y_{\tau-n:\tau},\,u_{\tau-n:\tau}\bigr),
 \label{eq:aug_encoder}
\end{equation}
where $\tau$ is the window start, $n$ the length of the encoder history,
$y_{\tau-n:\tau}$ stacks the measured positions $y_{\tau-n},\dots,y_\tau$,
and $u_{\tau-n:\tau}$ the recorded inputs at the same samples.
Both histories include the current sample $\tau$, whereas the objective
of~\cite[Eq.~(22e)]{hoekstra2026lfr} stops at $\tau-1$.
The reason is the initialisation below: the reconstructability map expresses
the state at $\tau$ through the outputs and inputs up to and including
$\tau$~\cite[Eq.~(15)]{hoekstra2026encoder}.
The encoder estimates all states at $\tau$, the measured positions included.

The encoder is a linear map with a nonlinear
correction~\cite[Eq.~(8)]{hoekstra2026encoder},
\begin{equation}
 \psi_\eta(\xi_\tau)
 =\begin{bmatrix}W^b\\ W^a\end{bmatrix}\xi_\tau+\tilde\psi(\xi_\tau),
 \label{eq:aug_enc_lin}
\end{equation}
where $W^b$ and $W^a$ give the physical and the additional states,
$\tilde\psi$ is an MLP, and $\eta$ collects $W^b$, $W^a$ and the parameters
of $\tilde\psi$.
The physical rows $W^b$ are initialised from the baseline through its
reconstructability map.
That map is derived for an LTI model, so we linearise \eqref{eq:eom} at the
nominal parameters about rest at $Y=0$, an equilibrium of the baseline, as
proposed for nonlinear baselines
in~\cite[Eqs.~(30), (31)]{hoekstra2026encoder}.
This gives
\begin{equation}
 A_c(Y)=\begin{bmatrix}0 & I\\ -M(Y)^{-1}K & -M(Y)^{-1}C\end{bmatrix},\qquad
 B_c(Y)=\begin{bmatrix}0\\ M(Y)^{-1}\end{bmatrix},
\label{eq:aug_lin}
\end{equation}
here evaluated at $Y=0$.
The zero-order-hold discretisation $(A_d,B_d)$ of $(A_c(0),B_c(0))$ over
$T_s$ matches the input held between samples, and
$C_d=\begin{bmatrix}P^\top & 0\end{bmatrix}$.
For this model, $W^b$ is the reconstructability
map~\cite[Eqs.~(16), (17)]{hoekstra2026encoder}
\begin{equation}
 W^b=\begin{bmatrix}
  A_d^nO_n^{+} & \;R_n-A_d^nO_n^{+}T_n
 \end{bmatrix},
 \label{eq:aug_wb}
\end{equation}
where $O_n$ is the observability matrix of $(A_d,C_d)$ over the history
length $n$, $T_n$ the Toeplitz matrix of the input-to-output responses,
$R_n$ the input-to-state map over the same history, and the two column blocks
act on the output and the input parts of $\xi_\tau$.
All matrices are scaled as the signals.
The measured positions make the linear model observable, so $O_n$ has full
column rank and its pseudoinverse $O_n^{+}$ is a left inverse.
The linearisation only sets the starting value of $W^b$: the encoder is
trained jointly with the model and can adapt to the nonlinear dynamics and to
other payload positions.

The rows $W^a$ have no baseline counterpart and are drawn randomly, and the
output layer of $\tilde\psi$ starts at zero.
At initialisation, the encoder therefore returns the reconstructed physical
state and a random additional state.
By \eqref{eq:aug_zero_init}, this additional state does not affect the
output, and for the reason given in Section~\ref{sec:aug_structure}, $W^a$
receives a gradient only once $f_{\mathrm{aug}}$ depends on $\bar x$.

\subsection{Closed-loop rollout and objective}\label{sec:aug_closed_loop}
The records are measured under feedback, and the model is used to predict the
machine under feedback: the response to a given reference, feedforward and
controller, for example to evaluate a changed controller before it is
deployed.
The truncated objective of~\cite[Eq.~(22)]{hoekstra2026lfr} instead
simulates each window open loop, driven by the recorded input, and
therefore scores a different response from the one the model is used for.
The baseline gives a second reason to leave that form.
$X$ and $Y$ have no stiffness in \eqref{eq:damping_stiffness_matrices}, so
$A_c(Y)$ in \eqref{eq:aug_lin} has two eigenvalues at the origin for every
$Y$.
In an open-loop rollout, a velocity error on these axes integrates into a
position error that persists over the rest of the window.
Inside the loop, the controller acts on such errors as it does on the
machine, provided the model in feedback with the controller is stable.

Kessels trains an augmented model with the known feedback controller inside
each truncated window and propagates the gradients through the
controller~\cite[Eqs.~(5.12), (5.13c), (5.13d)]{kessels2025ai}.
We adopt this rollout and keep the controller outside the learned model, as a
fixed and known part of the loop, so that it can be replaced in evaluation.
Kessels evaluates his model in this way with controller gains that differ
from those used in training~\cite[Sec.~5.3.2.3]{kessels2025ai}.

The recorded loop is driven by a reference $r$ and a feedforward force
$u^{\mathrm{ff}}$ through a linear time-invariant controller with state $s$
and matrices $(A_{\mathrm{fb}},B_{\mathrm{fb}},C_{\mathrm{fb}},D_{\mathrm{fb}})$,
which maps the position error to actuator forces.
The model loop receives the same $r$, $u^{\mathrm{ff}}$ and controller, and
differs only in its output $\hat y$ (Fig.~\ref{fig:aug_closed_loop}):
\begin{equation}
\begin{aligned}
 s_{k+1}&=A_{\mathrm{fb}} s_k+B_{\mathrm{fb}}(r_k-y_k),\\
 u_{\mathrm{act},k}&=C_{\mathrm{fb}} s_k+D_{\mathrm{fb}}(r_k-y_k)
   +u^{\mathrm{ff}}_k,\\
 \hat s_{k+1}&=A_{\mathrm{fb}} \hat s_k+B_{\mathrm{fb}}(r_k-\hat y_k),\\
 \hat u_{\mathrm{act},k}&=C_{\mathrm{fb}} \hat s_k
   +D_{\mathrm{fb}}(r_k-\hat y_k)+u^{\mathrm{ff}}_k ,
\end{aligned}
\label{eq:aug_direct_controller}
\end{equation}
where the hat marks the signals of the model loop, whose model input is
$\hat u_k=P\,\hat u_{\mathrm{act},k}$.
Subtracting the recorded loop from the model loop removes $r$ and
$u^{\mathrm{ff}}$,
\begin{equation}
\begin{aligned}
 x^c_{k+1}&=A_{\mathrm{fb}} x^c_k+B_{\mathrm{fb}}(y_k-\hat y_k),
   \qquad x^c_\tau=0,\\
 \hat u_{\mathrm{act},k}&=u_{\mathrm{act},k}+C_{\mathrm{fb}} x^c_k
   +D_{\mathrm{fb}}(y_k-\hat y_k),
\end{aligned}
\label{eq:aug_controller}
\end{equation}
where $x^c=\hat s-s$ is the difference of the controller states.
The model thus receives the recorded input plus the controller's response to
its own output error, and the rollout needs neither the reference nor the
feedforward.
The subtraction is exact when both loops apply the same linear controller,
not scheduled on the model state, with the same feedforward and without
actuator saturation.
When the controller of the model loop differs from the recorded one, the
subtraction no longer removes $r$ and $u^{\mathrm{ff}}$, and the model loop
has to be closed in the direct form \eqref{eq:aug_direct_controller}.

\begin{figure}[t]
 \centering
 \includegraphics[width=\columnwidth]{closed_loop_same_reference.pdf}
 \caption{Recorded plant loop (top) and augmented-model loop (bottom),
 driven by one shared reference $r_k$ and feedforward $u_k^{\mathrm{ff}}$,
 with the same feedback controller. The model output is evaluated
 before the input advances its state to the next sample.}
 \label{fig:aug_closed_loop}
\end{figure}

The condition $x^c_\tau=0$ starts the model's controller in the state of the
recorded controller, $\hat s_\tau=s_\tau$, which is the controller state the
machine had at $\tau$, so model error before $\tau$ is not carried into the
window.
The value of $s_\tau$ itself is never needed, because the recorded
$u_{\mathrm{act}}$ already contains its effect, and with $x^c_\tau=0$ a
model whose output equals the measurement receives exactly the recorded
input.
Kessels starts each window in the same controller state, but obtains it by
reconstructing the machine's controller state from the measured output, the
reference and the controller~\cite[Remark~5.4]{kessels2025ai}; the residual
form \eqref{eq:aug_controller} makes this reconstruction unnecessary.
Because the model has no feedthrough while the controller has
($D_{\mathrm{fb}}\neq0$), each sample is evaluated in a fixed order:
$\hat y_k$ from $\tilde x_k$, then $\hat u_{\mathrm{act},k}$ from
\eqref{eq:aug_controller}, then the model and controller states advance to
$k+1$.

Over a set $\mathcal T$ of window starts and a window length $n_f$, the
parameters $\vartheta=(\theta_{\mathrm{base}},\theta_{\mathrm{aug}},\eta)$
minimise the output error
\begin{equation}
 V(\vartheta)=\frac{1}{3\,|\mathcal T|\,n_f}
 \sum_{\tau\in\mathcal T}\sum_{\ell=0}^{n_f-1}
 \norm{y_{\tau+\ell}-\hat y_{\tau+\ell|\tau}}_2^2
 \label{eq:aug_loss}
\end{equation}
subject to \eqref{eq:aug_transition} with input $\hat u_k$,
\eqref{eq:aug_encoder} and \eqref{eq:aug_controller}, where
$\hat y_{\tau+\ell|\tau}$ is the prediction $\ell$ samples after the window
start, both outputs are normalised, and the factor $3$ is the number of
measured positions.
This is the truncated objective of~\cite[Eq.~(22)]{hoekstra2026lfr} with each
window simulated in closed loop as in~\cite[Eq.~(5.12)]{kessels2025ai}.
The physical combinations, the network and the encoder are estimated jointly,
while the controller is fixed.
No parameter regularisation term such as~\cite[Eq.~(26)]{hoekstra2026lfr} is
added; the split between the components is addressed by the construction of
Section~\ref{sec:obc} instead.
Every sample of the window is scored, from $\ell=0$.
Three horizons enter the method and are kept apart: the encoder history $n$,
the window length $n_f$ over which each rollout is scored, and the
validation rollout, which runs the same closed loop over each complete
validation record from one encoded state and selects the checkpoint.
Their values are given in Section~\ref{sec:setup}.
